\thispagestyle{empty}
\section*{Zusammenfassung}
In der vorliegenden Arbeit, HSP 2, wird untersucht, wie Straßenschäden mittels künstlicher Intelligenz genau klassifiziert werden können. Da konventionelle Methoden Defizite bei der Erfassung feiner Details aufweisen, finden Transformer-Modelle eine bessere Anwendung. Diese Modelle ermöglichen durch Aufmerksamkeits-Mechanismen ein besseres Verständnis der Schadensstrukturen.

Es erfolgt ein Vergleich zwischen dem klassischen Vision-Transformer, und dem hierarchischen Swin Transformer. Als Datenbasis wird der internationale Datensatz RDD2022 genutzt, der eine breite Vielfalt an Schadensbildern aus verschiedenen Ländern enthält. Der Datensatz wird zunächst vorverarbeitet und im Anschluss auf Grafikkarten trainiert. Die erzielten Ergebnisse werden anhand von Metriken wie Genauigkeit, Präzision und Recall analysiert. Es wird gezeigt, dass die KI in der Lage ist, komplexe Schadensmuster zuverlässig zu erkennen und von natürlichen Asphaltstrukturen zu unterscheiden. 

Die erzielten Resultate zeigen eine hohe Leistungsfähigkeit mit einer Klassifikationsgenauigkeit von bis zu 96,4\,\%. Insbesondere der Swin Transformer weißt eine bessere Erkennung Straßenschäden auf, insbesondere auf den Aspekt der Performance.

Abschließend wird aufgezeigt, wie diese Technologie in der Praxis unter Berücksichtigung nationaler Normen, beispielsweise ZTV ZEB, eingesetzt werden kann. Es wird dargelegt, dass der Einsatz künstlicher Intelligenz eine schnellere, objektivere und kosteneffizientere Sanierungsplanung in Bauämtern unterstützt.
\newpage
